\documentclass[12pt]{article}
\usepackage[margin=1in]{geometry}
\usepackage{amsmath}
\usepackage{amssymb}

\begin{document}

\section*{ECO 310 - Problem Set 2}
\subsection*{Question 5(c) - Full Solution}

\textbf{Problem:} Assume that $u(w) = 1 - e^{-\alpha w}$, where $\alpha > 0$ (CARA utility function). Solve explicitly for the optimal insurance level $q^*$, and show that it does not depend on wealth $w_0$.

\vspace{0.5cm}

\textbf{Solution:}

\textbf{Step 1: Find the marginal utility}

Given $u(w) = 1 - e^{-\alpha w}$, we have:
\begin{equation}
u'(w) = \alpha e^{-\alpha w}
\end{equation}

\vspace{0.3cm}

\textbf{Step 2: Apply the first-order condition from part (b)}

From part (b), the optimal insurance level satisfies:
\begin{equation}
\frac{u'(w_0 - L + (1-p)q)}{u'(w_0 - pq)} = \frac{p(1-\pi)}{\pi(1-p)}
\end{equation}

\vspace{0.3cm}

\textbf{Step 3: Substitute the CARA marginal utility}

Substituting $u'(w) = \alpha e^{-\alpha w}$ into both numerator and denominator:
\begin{equation}
\frac{\alpha e^{-\alpha(w_0 - L + (1-p)q)}}{\alpha e^{-\alpha(w_0 - pq)}} = \frac{p(1-\pi)}{\pi(1-p)}
\end{equation}

The $\alpha$ terms cancel:
\begin{equation}
\frac{e^{-\alpha(w_0 - L + (1-p)q)}}{e^{-\alpha(w_0 - pq)}} = \frac{p(1-\pi)}{\pi(1-p)}
\end{equation}

\vspace{0.3cm}

\textbf{Step 4: Simplify using properties of exponents}

Using the rule $\frac{e^a}{e^b} = e^{a-b}$:
\begin{equation}
e^{-\alpha(w_0 - L + (1-p)q) - (-\alpha(w_0 - pq))} = \frac{p(1-\pi)}{\pi(1-p)}
\end{equation}

Expand the exponent:
\begin{align}
&-\alpha(w_0 - L + (1-p)q) + \alpha(w_0 - pq) \\
&= -\alpha w_0 + \alpha L - \alpha(1-p)q + \alpha w_0 - \alpha pq \\
&= \alpha L - \alpha(1-p)q - \alpha pq \\
&= \alpha L - \alpha q[(1-p) + p] \\
&= \alpha L - \alpha q
\end{align}

Notice that \textbf{$w_0$ cancels out completely!}

Therefore:
\begin{equation}
e^{\alpha L - \alpha q} = \frac{p(1-\pi)}{\pi(1-p)}
\end{equation}

\vspace{0.3cm}

\textbf{Step 5: Solve for $q^*$}

Taking the natural logarithm of both sides:
\begin{equation}
\alpha L - \alpha q = \ln\left(\frac{p(1-\pi)}{\pi(1-p)}\right)
\end{equation}

Solving for $q$:
\begin{equation}
\alpha q = \alpha L - \ln\left(\frac{p(1-\pi)}{\pi(1-p)}\right)
\end{equation}

\begin{equation}
\boxed{q^* = L - \frac{1}{\alpha}\ln\left(\frac{p(1-\pi)}{\pi(1-p)}\right)}
\end{equation}

\vspace{0.5cm}

\textbf{Conclusion:}

The optimal insurance level $q^*$ depends only on the loss amount $L$, the risk aversion parameter $\alpha$, the probability of an accident $\pi$, and the insurance price $p$.

Crucially, \textbf{$q^*$ does not depend on initial wealth $w_0$}, which is the defining property of CARA (Constant Absolute Risk Aversion) preferences. The demand for insurance is independent of wealth level.

\end{document}
